\section{Segundo Parcial 2017}

% NOTA: el encabezado del PDF dice ``Parcial 2 -- 2015''; se conserva el rótulo
% asignado por la consigna (``Segundo Parcial 2017'').

\subsection{Ejercicio 1 — Modelo de compartimentos de confidencialidad (Bell--LaPadula)}

\begin{consigna}
Un famoso restaurante de 5 estrellas en la guía Michelin se destaca por todo lo que sirve: platos salados, platos dulces y bebidas (vinos y tragos).

Se tienen los siguientes sujetos:
\begin{itemize}
  \item Jorge, José y Juan son los dueños.
  \item Manuel y Ramón son chefs.
  \item Luis es uno de los mozos.
  \item Andrea es una cliente.
\end{itemize}
Considerar los siguientes objetos:
\begin{itemize}
  \item Recetas (tanto de platos dulces, como de salados, como de tragos).
  \item Menu. (se considera menú al listado de platos con sus respectivos precios).
  \item Listas de Stock (de productos para platos dulces, salados y tragos).
  \item Comanda. (pedido que se hace al camarero-mozo).
\end{itemize}
Diseñar un modelo de compartimentos de confidencialidad, que permita garantizar los siguientes requisitos, justificando el cumplimiento de cada uno.
\begin{itemize}
  \item Jorge puede leer todas las recetas, las listas de stock, el menú y las comandas.
  \item José, el somellier, es el único que escribe recetas de tragos. También lee stock de productos para tragos.
  \item Juan es el único que escribe recetas de platos (tanto dulces como salados). También lee stock de productos para platos.
  \item Los chefs pueden leer las recetas, pero no modificarlas.
  \item Los chefs pueden escribir en las listas de stocks para recomendar la compra de productos de su área.
  \item Ramón es el jefe de la barra, y sólo se ocupa de armar tragos. No cocina platos.
  \item Manuel dirige la cocina de platos dulces y salados.
  \item Los mozos sólo leen el menú y escriben comandas (no leen ni escriben nada mas).
  \item Los clientes sólo leen menús.
\end{itemize}
\end{consigna}

\paragraph{Temas.} Clase 6 — Políticas y Control de Acceso: Bell--LaPadula con
\emph{compartimentos}, relación de dominancia, reticulado, \emph{need-to-know}, reglas
\emph{read-down / write-up} (\S Compartimentos, reticulado y dominancia).

\paragraph{Qué necesitás saber.} Bell--LaPadula es un modelo de \textbf{confidencialidad}.
Un nivel de seguridad es un par $(L,C)$: una \textbf{clasificación} $L$ (orden total) y un
conjunto de \textbf{categorías/compartimentos} $C$. La relación de dominancia es
\begin{equation}
  (L,C) \succeq (L',C') \iff L' \le L \;\wedge\; C' \subseteq C .
\end{equation}
Las reglas de acceso:
\begin{itemize}
  \item \textbf{Seguridad simple (read-down):} $s$ lee $o$ si $L(s) \succeq L(o)$.
  \item \textbf{Propiedad $\ast$ (write-up):} $s$ escribe $o$ si $L(o) \succeq L(s)$.
\end{itemize}
Como aquí todos los requisitos pesan sobre \emph{quién accede a qué área} (tragos, platos
dulces, salados) y no sobre jerarquías de confianza, el problema se resuelve casi por
completo con \textbf{categorías} (need-to-know): un mismo nivel base y compartimentos por
área. Pares incomparables (ninguna $C$ es subconjunto de la otra) quedan
\textbf{efectivamente separados} (ni lectura ni escritura).

\paragraph{Notas y conexiones.} El enunciado pide ``compartimentos de confidencialidad'',
que es exactamente la extensión por categorías de Bell--LaPadula. Cuidado con la dirección:
para que un mozo \emph{escriba} una comanda que un dueño \emph{lee}, la comanda debe estar en
un nivel que el mozo domine al escribir (write-up) y el dueño domine al leer (read-down) — es
el patrón clásico ``el de abajo escribe hacia arriba, el de arriba lee hacia abajo''. Invertir
read-down/write-up es el error que más penaliza (Clase 6, \S Bell--LaPadula).

\paragraph{Resolución.} Usamos dos clasificaciones, \textbf{Alta} (dueños/dirección) $>$
\textbf{Baja} (resto), y categorías por área:
\[
  \mathcal{C} = \{\,\textsf{tragos},\ \textsf{dulces},\ \textsf{salados},\ \textsf{operación}\,\},
\]
donde \textsf{operación} agrupa lo público del salón (menú y comandas). Asignación de
niveles a \textbf{objetos}:
\begin{center}
\renewcommand{\arraystretch}{1.2}
\begin{tabular}{ll}
\toprule
Objeto & Nivel $(L,C)$\\
\midrule
Receta de tragos            & $(\text{Alta},\{\textsf{tragos}\})$\\
Receta dulces / salados     & $(\text{Alta},\{\textsf{dulces}\})$ / $(\text{Alta},\{\textsf{salados}\})$\\
Stock de tragos             & $(\text{Baja},\{\textsf{tragos}\})$\\
Stock dulces / salados      & $(\text{Baja},\{\textsf{dulces}\})$ / $(\text{Baja},\{\textsf{salados}\})$\\
Menú                        & $(\text{Baja},\{\textsf{operación}\})$\\
Comanda                     & $(\text{Baja},\{\textsf{operación}\})$\\
\bottomrule
\end{tabular}
\end{center}

Niveles de \textbf{sujetos} y verificación (recordando read-down para leer, write-up para
escribir):
\begin{itemize}
  \item \textbf{Jorge (dueño, supervisor):} $(\text{Alta},\{\textsf{tragos},\textsf{dulces},\textsf{salados},\textsf{operación}\})$.
        Domina todas las recetas, stocks, menú y comandas $\Rightarrow$ \emph{lee} todo
        (read-down). Cumple ``puede leer todas las recetas, stocks, menú y comandas''.
  \item \textbf{José (somellier, escribe recetas de tragos):} para \emph{escribir} la receta
        de tragos $(\text{Alta},\{\textsf{tragos}\})$ necesita $L(o)\succeq L(s)$, así que
        $L(s)=\text{Baja}$ o $\text{Alta}$ con $C \subseteq \{\textsf{tragos}\}$. Tomamos
        $(\text{Baja},\{\textsf{tragos}\})$: escribe receta de tragos (write-up, Alta domina a
        Baja) y \emph{lee} stock de tragos $(\text{Baja},\{\textsf{tragos}\})$ (mismo nivel,
        read-down). Es el único con categoría \textsf{tragos} y permiso de escritura de receta
        $\Rightarrow$ ``único que escribe recetas de tragos''.
  \item \textbf{Juan (escribe recetas de platos dulces y salados):}
        $(\text{Baja},\{\textsf{dulces},\textsf{salados}\})$: escribe ambas recetas (write-up) y
        lee ambos stocks. Único con esas categorías y escritura de receta.
  \item \textbf{Ramón (chef, jefe de barra — solo tragos):}
        $(\text{Baja},\{\textsf{tragos}\})$ pero \textbf{sin} derecho de escritura sobre la
        receta (la escritura de receta es exclusiva de José). Por dominancia \emph{podría}
        leer la receta de tragos (read-down: $(\text{Baja},\{\textsf{tragos}\})$ no domina
        $(\text{Alta},\{\textsf{tragos}\})$ \dots\ ver abajo) y \emph{escribir} en el stock de
        tragos (mismo nivel) $\Rightarrow$ ``escribe en stock para recomendar compras de su
        área''. No cocina platos: no tiene \textsf{dulces}/\textsf{salados}.
  \item \textbf{Manuel (chef, dirige cocina de dulces y salados):}
        $(\text{Baja},\{\textsf{dulces},\textsf{salados}\})$, lee recetas y escribe en los
        stocks de ambas áreas.
  \item \textbf{Luis y los mozos:} $(\text{Baja},\{\textsf{operación}\})$. \emph{Leen} menú
        (mismo nivel) y \emph{escriben} comandas. No tienen categorías de área $\Rightarrow$
        incomparables con recetas/stocks: ni leen ni escriben nada más.
  \item \textbf{Andrea y los clientes:} $(\text{Baja},\{\textsf{operación}\})$ con permiso de
        \emph{solo lectura} del menú (sin escritura de comanda, que es del mozo).
\end{itemize}

\textbf{Ajuste de niveles (lectura de recetas por chefs).} ``Los chefs pueden leer las recetas
pero no modificarlas'': leer exige $L(s)\succeq L(o)$. Para que Ramón/Manuel
\emph{lean} las recetas, las recetas deben estar a un nivel \emph{dominado por} los chefs, o
bien las recetas y los chefs comparten clasificación. La forma limpia es colocar las
\textbf{recetas en el mismo nivel base} que el escritor y darle al lector la misma categoría:
\begin{align*}
  \text{Receta tragos} &: (\text{Baja},\{\textsf{tragos}\}) \\
  \text{Ramón} &: (\text{Baja},\{\textsf{tragos}\}) \;\Rightarrow\; \text{lee (mismo nivel)},
\end{align*}
y la \textbf{exclusividad de escritura} de la receta (solo José/Juan) se garantiza con la
\textbf{ACL del objeto} (DAC), no con Bell--LaPadula: BLP dice quién \emph{puede} por flujo,
y la ACL restringe quién \emph{tiene} el permiso de escritura. Así un chef con
$(\text{Baja},\{\textsf{tragos}\})$ lee la receta (read-down/mismo nivel) pero no figura en la
ACL de escritura $\Rightarrow$ no la modifica. La separación efectiva entre áreas (un chef de
dulces no ve nada de tragos) la da que \textsf{dulces} y \textsf{tragos} son
\textbf{incomparables} en el reticulado.

\medskip
\textbf{Resumen del reticulado de categorías} (need-to-know): cada sujeto sólo recibe las
categorías de su área; Jorge las tiene todas (lee todo); mozos y clientes sólo
\textsf{operación}. La confidencialidad pedida se cumple porque los compartimentos
incomparables están separados y la escritura exclusiva se cierra con la ACL.

% ---------------------------------------------------------------------------
\subsection{Ejercicio 2 — Verdadero/Falso justificado (pentesting, Anderson, STRIDE)}

\begin{consigna}
Indica en cada caso si la afirmación es verdadera o falsa. Justificar en todos los casos. Si no se justifica correctamente, la respuesta quedará invalidada.
\begin{enumerate}[label=\alph*)]
  \item En las pruebas de pentesting siempre se intenta explotar las vulnerabilidades encontradas.
  \item Si un atacante puede probar $10^4$ passwords por segundo, y las passwords están compuestas por símbolos en [a-zA-Z0-9], entonces las passwords deberían tener 5 símbolos y 5 bits de SALT para que la probabilidad de adivinarlas en el lapso de un año sea menor a 0,5.
  \item STRIDE forma parte del método de Modelado de Amenazas de Microsoft.
\end{enumerate}
\end{consigna}

\paragraph{Temas.} (a) Clase 10b — Pentesting (metodología, explotar como último recurso).
(b) Clase 7 — Autenticación (fórmula de Anderson $P=TG/N$, salting). (c) Clase 8 —
Principios de Diseño y Vulnerabilidades (Threat Modeling de Microsoft, STRIDE).

\paragraph{Qué necesitás saber.}
\begin{itemize}
  \item \textbf{Pentesting:} explotar es el \textbf{último recurso}, no algo que se haga
        siempre; a menudo basta con analizar documentación o comportamiento, y la prueba debe
        ser lo \emph{menos intrusiva posible} (Clase 10b, \S Metodología de hipótesis de falla).
  \item \textbf{Anderson:} $P = \dfrac{T\cdot G}{N}$ con $P$ prob.\ de adivinar, $T$ tiempo,
        $G$ pruebas/seg, $N$ espacio de claves. El \textbf{salt no agranda} el espacio de
        claves $N$ que un atacante recorre para \emph{una} cuenta: sólo impide
        \textbf{paralelizar} el ataque entre cuentas. No entra en la cuenta de $P$ por cuenta.
  \item \textbf{STRIDE:} es el modelo de 6 categorías (Spoofing, Tampering, Repudiation,
        Information disclosure, DoS, Elevation of privilege) que Microsoft usa dentro de su
        Threat Modeling, cruzándolo con el DFD (Clase 8, \S STRIDE).
\end{itemize}

\paragraph{Notas y conexiones.} El ítem (b) es la trampa típica de Anderson: hay que
\emph{calcular} y comparar, y entender que ``bits de salt'' no afectan $P$ para una sola
cuenta. El alfabeto [a-zA-Z0-9] tiene $26+26+10 = 62$ símbolos.

\paragraph{Resolución.}

\textbf{(a) FALSO.} En pentesting \textbf{no} siempre se explota: explotar es el
\emph{último recurso} y el menos eficiente. Muchas veces alcanza con analizar documentación,
configuración o comportamiento del sistema para confirmar la vulnerabilidad; además el test
debe ser lo menos intrusivo posible (explotar puede provocar un DoS) y siempre repetible y
concluyente (Clase 10b — Pentesting, \S Metodología de hipótesis de falla).

\medskip
\textbf{(b) FALSO.} Con $N = 62^{L}$, $G = 10^{4}$/s y $T = 1$ año
$= 365\cdot24\cdot3600 = 3.1536\times10^{7}$ s, pedimos $P<0.5$:
\begin{equation}
  P = \frac{T\,G}{N} < 0.5
  \;\Longrightarrow\; N > \frac{T\,G}{0.5}
  = \frac{3.1536\times10^{7}\cdot 10^{4}}{0.5}
  = 6.31\times10^{11}.
\end{equation}
Con $L=5$: $N = 62^{5} = 9.16\times10^{8}$, que es \textbf{mucho menor} que $6.31\times10^{11}$,
luego
\begin{equation}
  P = \frac{3.1536\times10^{11}}{9.16\times10^{8}} \approx 344 \gg 0.5
\end{equation}
(es decir, se adivina con certeza en mucho menos de un año). La longitud mínima necesaria es
\begin{equation*}
  L \ge \log_{62}\!\bigl(6.31\times10^{11}\bigr) \approx 6.6 \;\Rightarrow\; L \ge 7,
\end{equation*}
no $5$. Además, los \textbf{bits de SALT no cambian} esta cuenta: el salt no aumenta el espacio
$N$ que el atacante recorre para una cuenta dada; sólo evita reutilizar el cómputo \emph{entre}
cuentas (no se paraleliza el ataque). Por ambos motivos la afirmación es falsa (Clase 7 —
Autenticación, \S Fórmula de Anderson y \S Salting).

\medskip
\textbf{(c) VERDADERO.} STRIDE es precisamente el modelo de clasificación de amenazas (6
categorías: Spoofing, Tampering, Repudiation, Information disclosure, Denial of service,
Elevation of privilege) que forma parte del método de \textbf{Threat Modeling} creado por
Microsoft, donde se cruza cada categoría con los elementos del DFD (Clase 8 — Principios de
Diseño y Vulnerabilidades, \S STRIDE).

% ---------------------------------------------------------------------------
\subsection{Ejercicio 3 — Secreto compartido de Shamir y entropía}

\begin{consigna}
Considerar el conjunto $C$ de sombras para un esquema de secreto compartido de Shamir $(k,n)$, donde $k=3$ y $n=4$ sobre $\mathbb{Z}_7$.
\[
  C = \{(1,6),(2,1),(3,0),(4,3)\}
\]
\begin{enumerate}[label=\alph*)]
  \item Obtener el secreto.
  \item Obtener el polinomio utilizado para obtener el conjunto de sombras.
  \item Expresar el significado de un esquema de secreto umbral $(3,n)$ utilizando el concepto de entropía.
\end{enumerate}
\end{consigna}

\paragraph{Temas.} Clase 8 — Principios de Diseño y Vulnerabilidades: \emph{secreto compartido
de Shamir en términos de entropía} (\S El secreto compartido de Shamir); reconstrucción por
interpolación polinómica (Primera Parte / criptografía). Clase 9 — Flujo de Información:
entropía de Shannon (\S Entropía como medida de información).

\paragraph{Qué necesitás saber.} En Shamir $(k,n)$ el secreto es $f(0)=a_0$ de un polinomio
de grado $k-1$ sobre un cuerpo finito. Con $k=3$ el polinomio es
$f(x)=a_0+a_1x+a_2x^2 \pmod 7$. Con $k$ sombras se reconstruye resolviendo el sistema lineal
(o por interpolación de Lagrange). La validez del esquema se expresa con \textbf{entropía}:
con menos de $k$ sombras la entropía del secreto \textbf{no cambia} (no hay flujo de
información); con $k$ o más, cae a $0$ (Clase 8, \S El secreto compartido de Shamir).

\paragraph{Notas y conexiones.} Esto conecta la criptografía de la Primera Parte
(reconstrucción polinómica) con el marco de \emph{flujo de información} de la Segunda:
``revelar el secreto'' equivale a $H(S)\to 0$. Toda la aritmética es módulo $7$.

\paragraph{Resolución.}

\textbf{(b) Polinomio.} Buscamos $f(x)=a_0+a_1x+a_2x^2 \bmod 7$ que pase por
$(1,6),(2,1),(3,0)$:
\begin{align*}
  a_0 + a_1 + a_2 &\equiv 6 \pmod 7 \\
  a_0 + 2a_1 + 4a_2 &\equiv 1 \pmod 7 \\
  a_0 + 3a_1 + 9a_2 \equiv a_0 + 3a_1 + 2a_2 &\equiv 0 \pmod 7
\end{align*}
Restando la 1.\textsuperscript{a} a la 2.\textsuperscript{a} y a la 3.\textsuperscript{a}:
\begin{align*}
  a_1 + 3a_2 &\equiv -5 \equiv 2 \pmod 7 \\
  2a_1 + a_2 &\equiv -6 \equiv 1 \pmod 7
\end{align*}
De la primera $a_1 \equiv 2 - 3a_2$. Sustituyendo en la segunda:
$2(2-3a_2)+a_2 = 4 - 5a_2 \equiv 1 \Rightarrow -5a_2 \equiv -3 \Rightarrow 5a_2 \equiv 3$.
Como $5^{-1}\equiv 3 \pmod 7$ (pues $5\cdot3=15\equiv1$), $a_2 \equiv 3\cdot3 = 9 \equiv 2$.
Entonces $a_1 \equiv 2 - 3\cdot2 = -4 \equiv 3$ y $a_0 \equiv 6 - 3 - 2 = 1$. Por lo tanto
\begin{equation}
  \boxed{\,f(x) = 1 + 3x + 2x^2 \pmod 7\,}.
\end{equation}
\emph{Verificación} con la 4.\textsuperscript{a} sombra $(4,3)$:
$f(4)=1+12+32=45\equiv 45-42 = 3 \pmod 7$ \checkmark.

\medskip
\textbf{(a) Secreto.} El secreto es el término independiente:
\begin{equation}
  S = f(0) = a_0 = \boxed{1}.
\end{equation}

\medskip
\textbf{(c) Significado vía entropía.} Un esquema umbral $(3,n)$ significa que se necesitan
\textbf{al menos $3$ sombras} para recuperar el secreto $S$, y que con \emph{menos} de $3$ no
se obtiene \textbf{ninguna} información sobre $S$. En términos de entropía, llamando $S_i$ a
las sombras:
\begin{align*}
  H(S \mid S_{i_1}) &= H(S) && \text{(1 sombra: la incertidumbre no baja)}\\
  H(S \mid S_{i_1}, S_{i_2}) &= H(S) && \text{(2 sombras: tampoco hay flujo)}\\
  H(S \mid S_{i_1}, S_{i_2}, S_{i_3}) &= 0 && \text{(3 sombras: $S$ queda determinado)}
\end{align*}
Es decir: con menos de $3$ sombras \textbf{no hay flujo de información} ($H(S)$ se mantiene
constante), y al alcanzar el umbral $H(S)\to 0$ (certeza). Esa es la definición de que el
esquema es válido/seguro (Clase 8 — Principios de Diseño y Vulnerabilidades, \S El secreto
compartido de Shamir; Clase 9 — Flujo de Información, \S Entropía).

% ---------------------------------------------------------------------------
\subsection{Ejercicio 4 — Multiple choice (resolución de conflictos en ACL y DMZ)}

\begin{consigna}
En cada caso selecciona la opción correcta. Justificar en todos los casos. Si no se justifica correctamente, la respuesta quedará invalidada.
\begin{enumerate}
  \item Si se tiene una lista de control de acceso $\mathrm{ACL}(o) = \{(\text{Manuel}, r), (\text{Directivos}, w), (\ast, w)\}$. Asumiendo que Manuel es Directivo,
    \begin{enumerate}[label=\alph*.]
      \item Manuel puede leer y escribir en $o$ si la política es first rule.
      \item Manuel puede leer y escribir en $o$ si la política es grant all.
      \item Manuel puede leer y escribir en $o$ si la política es override.
      \item Manuel puede leer y escribir en $o$ si la política es augment.
    \end{enumerate}
  \item En una red bien diseñada, la zona demilitarizada (DMZ) nunca debería contener:
    \begin{enumerate}[label=\alph*.]
      \item Servidores de email.
      \item Servidores de DNS.
      \item Servidores de subred de desarrollo.
      \item Servidores de Web Proxy.
    \end{enumerate}
\end{enumerate}
\end{consigna}

\paragraph{Temas.} (4.1) Clase 6 — Políticas y Control de Acceso: resolución de conflictos en
ACL (first rule, grant all, override/best-match, augment) (\S Conflictos en ACL, \S Derechos
por defecto). (4.2) Clase 10 — Seguridad en la Empresa: Defense-in-depth y DMZ (\S DMZ).

\paragraph{Qué necesitás saber.}
\begin{itemize}
  \item \textbf{Conflictos de ACL} cuando $s$ matchea varias entradas (Clase 6):
    \begin{itemize}
      \item \textbf{Permitir (unión):} otorga la unión de los permisos.
      \item \textbf{Grant All (intersección):} deniega si \emph{alguna} entrada deniega
            $\Rightarrow$ sólo los permisos que \emph{todas} otorgan.
      \item \textbf{First rule / first match:} aplica la \emph{primera} entrada que matchea.
      \item \textbf{Override / best match:} usa la entrada \emph{más específica}.
      \item \textbf{Augment:} parte del \emph{default} ($\ast$) y le \emph{agrega} las
            entradas explícitas.
    \end{itemize}
  \item \textbf{DMZ:} zona expuesta entre el firewall externo y el interno; aloja los
        servicios que deben ser accesibles desde Internet (web, correo, DNS, proxy). Lo
        \emph{interno} (LAN, bases de datos internas, entornos de desarrollo) \textbf{nunca}
        va en la DMZ (Clase 10, \S DMZ).
\end{itemize}

\paragraph{Notas y conexiones.} Manuel matchea las tres entradas: $(\text{Manuel},r)$ por
nombre, $(\text{Directivos},w)$ por grupo y $(\ast,w)$ por default. El resultado ``$r$ y $w$''
depende de la política. Esto es el ejercicio resuelto análogo de Clase 6 (\S Conflictos en
ACL).

\paragraph{Resolución.}

\textbf{(4.1) Correcta: d) augment.}
Manuel coincide con las tres entradas. Sus permisos según cada política:
\begin{center}
\renewcommand{\arraystretch}{1.25}
\begin{tabular}{>{\bfseries}l l l}
\toprule
Política & Cómo resuelve & Permisos de Manuel\\
\midrule
First rule & primera entrada $(\text{Manuel},r)$ & solo $r$ \;(no)\\
Grant all (intersección) & $r \cap w \cap w$ & $\varnothing$ \;(no)\\
Override / best match & entrada más específica $(\text{Manuel},r)$ & solo $r$ \;(no)\\
Augment & default $\ast{:}w$ $+$ específicas $r,w$ & $r$ y $w$ \;(si)\\
\bottomrule
\end{tabular}
\end{center}
Sólo con \textbf{augment} (default $+$ entradas explícitas: $w$ del $\ast$, más $r$ de Manuel y
$w$ de Directivos) Manuel obtiene \textbf{lectura y escritura}. Con first rule y con override
queda sólo $r$ (la entrada más específica/primera es $(\text{Manuel},r)$), y con grant all la
intersección de $\{r\},\{w\},\{w\}$ es vacía. Por lo tanto la opción correcta es \textbf{d)}
(Clase 6 — Políticas y Control de Acceso, \S Conflictos en ACL / \S Derechos por defecto).

\medskip
\textbf{(4.2) Correcta: c) servidores de subred de desarrollo.}
La DMZ es la zona \emph{expuesta} hacia Internet y aloja servicios públicos: servidores web,
de correo (email), de DNS y de Web Proxy son todos servicios de cara al exterior que sí van en
la DMZ. Una \textbf{subred de desarrollo} es un recurso \emph{interno} (código, entornos no
endurecidos, datos de prueba) que debe quedar detrás del firewall interno, en la LAN, nunca en
una zona expuesta. Exponerla violaría la segmentación de Defense-in-depth. Por eso la DMZ
nunca debería contenerla (Clase 10 — Seguridad en la Empresa, \S DMZ).
